\section{Dihedral Groups}
Consider the set of all symmetries of a regular $n$-gon. They form a group under composition of transformations. We call this a \textit{Dihedral Group}. Denoted as $D_n$ (or $D_{2n}$ in some texts). Notice $D_n$ has $2n$ elements. Let's level $n$ vertices from $1$ to $n$. Then for $1$ we have $n$ choices. For $2$ we have $2$ choices since vertex $2$ is adjacent to vertex $1$. And since symmetries are rigid motions, once we specify vertex $1$ and $2$, positions of the remaining vertices are determined automatically. Hence,
$$|D_n| = 2n.$$
\begin{figure}[H]
    \centering
    \def\svgwidth{0.6\columnwidth}
    \incfig{n-gon}
    \caption{A regular n-gon}
    \label{fig:n-gon}
\end{figure}
Now, for $D_3$ or $D_4$ we can explicitely exhibit all the symmtries and do computations, but for higher $n$ it becomes complicated. Therefore, viewing $D_n$ as an abstract group is convenient. Label a regular $n$-gon consecutively from $1$ to $n$ clockwise. Let $r$ be the rotation anti-clockwise about the origin by $\frac{2 \pi}{n}$ radian. Let $s$ be the reflection about the line of symmetry through vertex $1$ and the origin. Then, the following properties are true:
\begin{enumerate}
    \item \(1, r, r^2, \ldots, r^{n-1}\) are all distinct and \(r^n = 1\), so \(|r| = n\).
    \item \(s^2 = 2\).
    \item $r^is$ has to be a reflection, that is, \(r^is \neq r^j\) for any \(i\) and \(j\).
    \item \(sr^i \neq sr^j\), for all \(0 \leq i, j \leq n - 1\) with \(i \neq j\), so
    \[
    D_{n} = \{1, r, r^2, \ldots, r^{n-1}, s, sr, sr^2, \ldots, sr^{n-1}\}
    \]
    i.e., each element can be written \textit{uniquely} in the form \(sr^i\) for some \(k = 0\) or \(1\) and \(0 \leq i \leq n - 1\).
    \item \(rs = sr^{-1}\). Because \[(rs)^2 = e \implies rs \cdot rs = e \implies rs = s^{-1}r^{-1} \implies rs = sr^{-1}.\]
    \item By induction, \(r^i s = sr^{-i}\), for all \(0 \leq i < n\).
\end{enumerate}
Using \textbf{(1)}, \textbf{(2)}, \textbf{(3)}, we can compute any finite number of produts of \(D_{n}\) and represent it in the form \(sr^i\), where \(k = 0\) or \(1\) and \(0 \leq i \leq n - 1\). For example, if \(n = 12\),
\begin{eg}
$(sr^9)(sr^6) = s(r^9 s)r^6 = s(sr^{-9})r^6 = s^2 r^{-9+6} = r^{-3} = r^9.$
\end{eg}
Notice, the set $\{1, r, r^2, \ldots, r^{n-1}\}$ is an abelian subgroup of $D_n$.
\section{Generators, Relations \& Presentation of a Group}
All the elements of $D_n$ is generated by $r, s$. We call them generators of $D_n$. $r, s$ satisfy some equations which are called relations. In general, if some group $G$ is generated by some $S \subseteq G$ and there is some collection of relations, say $R_1, R_2, ..., R_m$ ($R_i$ is an equation in the elements from $S \cup \{1\}$) such that any relation among the elements of $S$ can be deduced from these, we shall call these generators and relations a \textit{presentation} of $G$ and write $$G = \langle S | R_1, R_2, ..., R_m \rangle.$$
One presentation for the dihedral group $D_n$ is $$D_n = \langle r,s | r^n = s^2 = e, rs=sr^{-1}\rangle.$$
If $G$ is a finite group then, for $x, y \in G, \langle x, y \rangle$ is the set of any possible combination of products of $x, y$. $\langle x, y \rangle$ forms a subgroup of $G$. We call it \textit{free group} generated by $x, y$.



